\chapter{Dense self-attention and the KV cache}
\label{sec:quadratic}

\section{Vaswani attention}

For sequence length $n$, hidden dimension per head $d$, and one attention head, queries, keys, and values
are matrices $Q,K,V \in \mathbb{R}^{n\times d}$ \cite{vaswani2017attention}. Attention weights are
\begin{equation}
  A = \operatorname{softmax}\!\left(\frac{Q K^\top}{\sqrt{d}}\right) \in \mathbb{R}^{n\times n},
\end{equation}
and the output is $A V$. Thus each head performs $\Theta(n^2 d)$ work to form $QK^\top$ and multiply by $V$
(or equivalent fused kernels).

\section{Why $n^2$ hurts}

The matrix $A$ has $n^2$ entries per head per layer. At $n = 10^6$, that is $10^{12}$ scalars per head per layer.
Materializing $A$ in high-precision floating point is prohibitive at scale.

\section{KV cache at inference}

During autoregressive decoding, keys and values for past tokens must be retained. For $L$ layers,
$H$ heads, head dimension $d$, and $n$ tokens processed so far, KV storage scales linearly in $n$
(per layer), while \emph{dense} attention still entails $\Theta(n^2)$ pairwise scoring work unless the
algorithm changes.

\paragraph{Order-of-magnitude KV memory (illustrative).}
With both $K$ and $V$ stored in fp16, a rough per-token KV footprint scales like
$2 \cdot L \cdot H \cdot d$ bytes per token (constants depend on packing and sharing schemes such as MQA/GQA).
At large $n$, this dominates GPU memory before compute dominates wall clock---the ``KV cache wall'' emphasized
in applied systems work.
